% 2026-11-12 / B-v41. Current proof, not a version-history appendix.
\section{一般点復元の基礎体降下と定量的生成}
\label{gcd:sec:descent}

本節では可換模型の\emph{幾何学的一般点での存在}を得た後の議論を整理する。
方程式の定義体と、圏同値の定義体は異なる。
前者は低rotation積の可積分性により元の体へ降下できるが、
後者について本稿が保証するのは有限拡大後の存在である。
可積分性の有限線形代数そのものは第\ref{sec:explicit-jacobian-recovery}節と
第\ref{far:sec:recovery}節の結果を用いる。

\subsection{次数別flat入力から従う一般ファイバーのAS正則性}

\begin{proposition}[一般点のambientに対する追加仮定の除去]
\label{gcd:prop:automatic-as}
$(R,\mathfrak m,k)$をnoetherian局所整域、$K=\operatorname{Frac}R$とする。
$A_R$は非負次数の中央$R$-代数、$(A_R)_0=R$、各$(A_R)_j$は有限自由で、
$A_R/\mathfrak mA_R\simeq k[x_1,\ldots,x_N]$とする。
このとき$A_K$は左右noetherianな二次Koszul AS正則代数であり、
大域次元とGorenstein parameterは$N$、Hilbert級数は$(1-z)^{-N}$である。
さらに$\Omega_R\in Z(A_R)_d$の剰余が非零なら、$\Omega_R$は左右正則で、
商の各次数は$R$-自由である。
\end{proposition}
\begin{proof}
補題\ref{lp:lem:noetherian}の次数別Krull交叉の議論を$A_R$に適用すると、
左右noetherian性を得る。中央集合$R\setminus\{0\}$による局所化で$A_K$もそうなる。
ここでnoetherian性を有限Hilbert級数だけから推論していない。

補題\ref{lp:lem:resolution}を、augmentation加群$R$と、
$A_R^e=A_R\otimes_RA_R^{\mathrm{op}}$上の対角$A_R$に適用する。
多項式環のKoszul分解と対角分解を持ち上げると、項がそれぞれ
$A_R(-j)^{\binom Nj}$、$A_R^e(-j)^{\binom Nj}$である長さ$N$の分解を得る。
持上げの各段で、剰余上の生成元を持ち上げた写像の全射性を各内部次数の
Nakayamaで示す。有限自由$R$-加群への全射は分裂するので核も各次数で自由であり、
最後の核は剰余が零のため零となる。この構成は任意の通常の底変更に可換である。
したがって$A_K$は二次Koszulで、対角分解から左右大域次元は高々$N$となる。

augmentation分解の$A_R$双対を固定内部次数ごとに取る。
その剰余複体は、内部次数$-N$・コホモロジー次数$N$以外で完全である。
有界有限自由複体の単元行列成分を順に消去すれば、他の内部次数の複体は分裂完全となる。
内部次数$-N$では最後の$R$のみが残る。よって左右で
\begin{equation}\label{gcd:eq:as}
 R\underline{\operatorname{Hom}}_{A_R}(R,A_R)\simeq R(N)[-N]
\end{equation}
となり、底変更後の非零な最上Extから大域次元は正確に$N$である。
次数別自由性はHilbert級数を与え、増大度も多項式的である。
最後に各次数の$\Omega_R$乗法は剰余上で単射である。
列数と同じ大きさの単元小行列式を使えば$R$上でも分裂単射となり、
余核も自由となる。中央性により左右で同じ結論を得る。
\end{proof}

\subsection{smooth properなpairから取り出す低次数積}

$k$を標数零の体、$\mathcal C$をsmooth properな$k$-線形dg圏、
$\rho$をdg Morita自己同値とする。対角を$\Delta$、$\rho$の双加群を$P$と書き、
\begin{equation}\label{gcd:eq:low}
 V=H^0R\operatorname{Hom}_{\mathcal C^e}(\Delta,P),\qquad
 U=H^0R\operatorname{Hom}_{\mathcal C^e}(\Delta,P\otimes_{\mathcal C}^{\mathbf L}P)
\end{equation}
と置く。積$\mu:V\otimes V\to U$は$\rho(\beta)\circ\alpha$である。
通常の三角圏上の自然変換と、この導来自然変換を区別する。

\begin{lemma}[低次数積の体拡大との両立]\label{gcd:lem:basechange}
$V,U$は有限次元であり、任意の体拡大$\ell/k$に対して、その形成と積は底変更と可換である。
\end{lemma}
\begin{proof}
smoothnessにより$\Delta$は$\mathcal C^e$上perfectである。
proper性と$P$の可逆性により、$P$と$P\otimes_{\mathcal C}^{\mathbf L}P$は
基礎体上perfectである。有限cellとretractで$\Delta$を表示すると、
式\eqref{gcd:eq:low}のHom複体も$k$-perfectとなる。
同じ有限表示を使えばHomと平坦な体拡大が可換であり、$H^0$も可換である。
双加群のテンソルと射の合成はこの底変更に自然なので、積も保たれる。
\end{proof}

\begin{theorem}[可換cubic模型の方程式は基礎体上に取れる]
\label{gcd:thm:field-model}
$n\ge3$、$N=n+2$とする。上のpairがある代数閉拡大体$\Omega/k$上で
滑らかなcubic $Y/\Omega$のrotation付きKuznetsov pairとMorita同値であると仮定する。
このとき次を満たす$G\in\operatorname{Sym}^3_k V$が存在する。
\begin{enumerate}
\item $X_G=V(G)\subset\mathbb P(V^*)$は$k$上幾何学的に滑らかで、$X_{G,\Omega}\simeq Y$である。
\item ある有限拡大$k'/k$上で$(\mathcal C_{k'},\rho)\simeq(\operatorname{Ku}(X_{G,k'}),\rho_G)$である。
元の同値に三次周期比較を指定していた場合、それも保存できる。
\item $\mathcal C$の任意の滑らかなcubic模型は、共通の代数閉拡大体上で射影同型である。
\end{enumerate}
従って方程式の定義体には有限拡大を必要としない。
一方、同値双加群の$k$への降下や、$k$上での模型の一意性は主張しない。
\end{theorem}
\begin{proof}
命題\ref{prop:ej-geometric-low-degree}に用いる
Lin--Zhangのsector計算\cite[Theorem~3.5, Remark~3.6]{LZ}を$\Omega$上で適用する。
三次式の非恒等sectorがコホモロジー次数零へ寄与する内部次数は、
$N$が偶数の場合の$N/2$のみである。$N\ge5$では内部次数$1,2$に干渉しない。
したがって、$\Omega$上で$V,U,\mu$はJacobian環の次数$1,2$と積に一致する。
この計算は標数零の代数閉体上の代数的なfactorization計算である。
補題\ref{gcd:lem:basechange}と忠実平坦性から、$k$上でも
\begin{equation}\label{gcd:eq:gradient-data}
 \dim V=N,\quad \dim U=\binom N2,\quad
 \overline\mu:\operatorname{Sym}^2V\twoheadrightarrow U,\quad
 Q:=\ker\overline\mu,\quad\dim Q=N
\end{equation}
となる。積の対称性も降下する。
$Q$には代数閉包上で共通射影零点がない。これは$\Omega$上で勾配空間となることと
$Y$の滑らかさから従う。

$V$の基底を$x_i$、$Q$の基底を$q_a$とする。可積分性の線形解空間
\begin{equation}\label{gcd:eq:integrability}
 \mathcal L_Q=\left\{A\in\operatorname{Mat}_N(k):
 \sum_a A_{ia}\partial_jq_a=\sum_a A_{ja}\partial_iq_a\quad(i<j)\right\}
\end{equation}
は体拡大と可換である。$Y$の方程式の勾配を$q_a$で展開すると、
$\mathcal L_Q\otimes_k\Omega$に可逆行列が存在する。
$k$-基底$A_1,\ldots,A_r$を選べば
$D(z)=\det(\sum_\nu z_\nu A_\nu)$は零多項式でない。
$k$は無限体なので$D(c)\ne0$となる$c\in k^r$が存在する。
実際、各変数の次数は高々$N$なので$\{0,1,\ldots,N\}^r$のいずれかで非零となる。
$A=\sum c_\nu A_\nu$として
\begin{equation}\label{gcd:eq:model}
 G=\frac13\sum_{i,a}x_iA_{ia}q_a
\end{equation}
と置く。可積分性とEuler等式から$\partial_iG=\sum_aA_{ia}q_a$である。
従って$G$の勾配空間は$Q$、その共通射影零点はなく、$X_G$は幾何学的に滑らかである。

$\Omega$上で$G$と$Y$の方程式は同じ勾配空間を持つ。
命題\ref{prop:ej-projective-uniqueness}の対称三重線形形式の証明は
任意の標数零代数閉体で通るので、両者は線形座標変換で一致する。
この証明はsymmetrizer代数の可換性と非零冪零元の不存在を用い、
固有空間ごとの非零係数を三乗根で吸収するものである。
その座標変換はgraded MFの内部shiftと三次周期性を保つ。
従って$\Omega$上で$(\mathcal C,\rho)$と$X_G$のpairの同値を得る。

$X_G$は既に$k$上にあり、そのKuznetsov成分はsmooth properである。
固定された二つのproper dg模型に補題\ref{rog:lem:descent}を適用する。
perfect同値核、rotation比較、指定周期図式を有限生成$k$-部分代数へ降ろし、
単位・余単位のconeを消す非空開の閉点を選べば、有限拡大$k'/k$上の同値を得る。
第三の主張も、任意の模型に式\eqref{gcd:eq:gradient-data}を適用し、
同じ勾配空間の射影的一意性を用いれば従う。
\end{proof}

\begin{remark}[存在証明と実計算の違い]\label{gcd:rem:limits}
この定理は幾何学的な可換模型の存在を\emph{仮定して}、その方程式の定義体を改善する。
任意の非可換圏で次元と積分条件を検査するだけでは圏同値は得られない。
また、$V,U,\mu$の有限性の証明は、その非可換入力での具体行列を計算したという意味ではない。
$k=\mathbb C(h)$のときも$G$の係数が$h=0$で正則、あるいは$G_0=F$とは限らない。
特殊点を含む代数化には第\ref{far:sec:recovery}節の局所入力を別途確認する必要がある。
\end{remark}

\begin{example}[同じ勾配空間でも基礎体上の同型類は一意でない]
\label{gcd:ex:forms}
$k=\mathbb C(h)$、$N\ge5$とし、
\[
 F=x_1^3+\cdots+x_N^3,\qquad
 H=h x_1^3+x_2^3+\cdots+x_N^3
\]
を考える。両方とも幾何学的に滑らかで、勾配空間は
$\langle x_1^2,\ldots,x_N^2\rangle$である。
$k(h^{1/3})$上では同型だが、$k$上では射影同型でない。
実際、両方のHessianの零因子は座標超平面の和である。
射影線形同型はこの超平面配置を保つのでmonomial行列で表される。
係数を比較すると$h$が$k$の三乗であることを要するが、
$h$進付値が1なのでこれは不可能である。
したがって、勾配空間から得る方程式は\emph{幾何学的な}同型類を定めるのであり、
基礎体上のformや同値双加群の降下障害を自動的に消すものではない。
\end{example}

\subsection{既知の生成定理による最短経路と一様な上界}

\begin{proposition}[三shiftの定量的生成：BFKの帰結]
\label{gcd:prop:bound}
$k$は標数零の代数閉体、$Y\subset\mathbb P^{N-1}_k$は滑らかなcubic、$N\ge5$とする。
任意の非零$E\in\operatorname{Ku}(Y)$について
\begin{equation}\label{gcd:eq:bound}
 \operatorname U\bigl(E\oplus\rho_YE\oplus\rho_Y^2E\bigr)
 \le 2N(N+1)-1
\end{equation}
となる。ここで$\operatorname U$は、有限直和・shift・直和因子を段数零として、
圏の全対象を得るためのcone段数の最小値である。
\end{proposition}
\begin{proof}
Orlov比較により$\rho_Y$はgraded singularity categoryの内部shiftとなる。
Ballard--Favero--Katzarkov\cite[Theorem~5.15]{GCDBFK}は、
$N$変数の次数$d$の孤立斉次超曲面について、任意の非零対象の$d$個のshiftの直和が
生成時間$2N\bigl(N(d-2)+1\bigr)-1$以下のstrong generatorとなることを証明している。
$d=3$を代入すればよい。滑らかな射影cubicのアフィン錐の特異点は原点だけである。
\end{proof}

上の生成性と上界はBFKの既知定理の帰結である。
第\ref{rog:sec:generic}節にはStevenson分類と次数再付与を用いる別証明も与える。
一般点復元の主証明では、上記のBFK定理を直接適用する。
BFKの論文は基礎体を代数閉・標数零としているので、
$\mathbb C((h))$へそのまま適用せず、まずその代数閉包へ底変更する。
一般点同値の有限拡大への降下はその後で行う。
この上界は同値核の行列サイズや必要な体拡大次数の上界ではない。

\begin{proposition}[固定三対象の比較による独立の十分条件]
\label{gcd:prop:three-object}
主定理の代数的入力および局所比較(L)を仮定する。
完備MF圏全体の比較(CF)を使用せず、以下の固定三対象の比較を与えるだけでも
定理\ref{gdr:thm:generic}の結論が成立する。
ある体拡大$\ell/K$、滑らかなcubic $Y/\ell$、非零$E\in\operatorname{Ku}(Y)$について、
\[
 \Gamma_K\otimes_K\ell\simeq
 R\operatorname{End}(E\oplus\rho_Y E\oplus\rho_Y^2 E)
\]
というenhancedなMorita比較があり、これが生成部分圏上のrotationと指定周期比較を保つ。
ここでは単に両側のcohomology群や低次数の積が一致するだけでは不十分である。
\end{proposition}
\begin{proof}
右辺のperfect加群圏は三対象のthick閉包である。
$\ell$の代数閉包上で命題\ref{gcd:prop:bound}を適用すれば、
それは$\operatorname{Ku}(Y)$全体となる。
以後は主定理の有限降下、coherent飽和、基礎体上の方程式抽出と同じである。
\end{proof}

この条件は既に構成された代数的sourceに対する\emph{固定三対象の実際の比較}であり、
全MF圏上の関手を先に構成することを要しない。
この命題は、主定理に用いた全MF比較よりも弱い十分条件を分離したものである。
他方、相対的な変形理論では全MF比較そのものが必要になるため、
その構成は第\ref{cft:sec:transport}節に詳述する。

\subsection{多変数局所基底への適用と接着の問題}

\begin{proposition}[任意次元の正則局所基底の一般点]
\label{gcd:prop:multibase}
$(R,\mathfrak m,\mathbb C)$を有限生成複素代数の正則局所環とし、
$K=\operatorname{Frac}R$とする。次数別有限自由な中心三次対$(A_R,\Omega_R)$の
特殊ファイバーを多項式環と滑らかなcubicとする。
条件(L)をこの$R$について、条件(CF)を$\mathfrak m$進完備化$\widehat R$について仮定する。
このとき定理\ref{gdr:thm:generic}の結論は$K$をこの分数体に置き換えて成立する。
\end{proposition}
\begin{proof}
正則局所環の完備化は整域で、$K\hookrightarrow L=\operatorname{Frac}\widehat R$である。
(CF)を固定三対象に制限してから$L$、次いでその代数閉包へ底変更する。
(L)のperfect Homにより特殊点と一般点のEulerベクトルは同じで、
そのノルムは$2\cdot3^{N-1}>0$である。
命題\ref{gcd:prop:bound}により像は全可換MF圏を生成する。
以後のperfect核の有限降下、有限拡大からのsmoothness降下、coherent飽和は
基礎体が$\mathbb C(h)$であることを使わない。
命題\ref{gcd:prop:automatic-as}と定理\ref{gcd:thm:field-model}で結論を得る。
\end{proof}

\begin{lemma}[多変数でも成立する形式Poisson消去の代数部分]
\label{gcd:lem:multi-koszul}
$\widehat R=\mathbb C[[t_1,\ldots,t_r]]$、$S=\mathbb C[x_1,\ldots,x_N]$とし、
$W\in\widehat R\otimes S_3$の剰余が滑らかな射影cubicを定めるとする。
$\pi\in\mathfrak m(\widehat R\otimes S_2\otimes\bigwedge^2\langle\partial_x\rangle)$が
$[\pi,W]=[\pi,\pi]=0$を満たせば、一意な
$\Xi\in\mathfrak m\widehat R\otimes\bigwedge^3\langle\partial_x\rangle$があり、
$\pi=[\Xi,W]$、$[\Xi,\pi]=0$である。
\end{lemma}
\begin{proof}
勾配Koszul複体は特殊ファイバーで正の外積次数において完全である。
固定した重み（係数次数と外積次数の二倍の和）ごとに有限自由複体を取り、
$\mathfrak m$進の次数ごとにcycleを持ち上げれば、完備化上でも同じ完全性を得る。
二次係数の二重vectorの原像は係数次数零の三重vectorである。
その差の原像は係数次数$-2$を要求するので一意である。
Jacobi恒等式により$[\Xi,\pi]$は係数次数一の四重vectorのKoszul cycleであり、
その原像は係数次数$-1$となるから零である。
\end{proof}

この補題から、周期変数$u$について
$\exp(\operatorname{ad}(\Xi/u))(-uW+\pi)=-uW$が従う。
実際の完備MF比較と相対windowの構成には、それぞれ
定理\ref{cft:thm:complete}と定理\ref{tw:thm:window}を用いる。
これらの局所的な比較をモジュライの異なるチャートにわたって接着し、
共通の相対模型を構成する問題は本節の結論に含めない。
